\section{Introduction}\label{sec:intro}

Neutral-atom quantum processors and simulators start every experimental cycle from a stochastically loaded array of optical tweezers, in which about half of the sites hold an atom. The atoms are sorted into a defect-free target pattern before the computation or the simulation begins. Defect-free arrays of hundreds of atoms support the programmable simulators and the logical-qubit processors that now operate~\cite{ebadi,bluvstein}, arrays of about 6100 atoms have been assembled~\cite{manetsch}, and an architecture that keeps more than 3000 atoms for hours reloads the array continuously~\cite{chiu}. The time and the loss probability of the sorting step enter the yield of every cycle, so the way the traps are moved is a design question of the experiment and not a detail of its control software.

Two families of machines sort atoms. Acousto-optic deflectors move one row or one column of traps at a time, so that the number of moves grows with the size of the array~\cite{endres,barredo}. A spatial light modulator (SLM) displays a hologram that forms all the traps, and a sequence of holograms moves every trap in parallel; the number of moves does not depend on the number of atoms. The SLM route was demonstrated with traps whose movement was designed analytically as a linear combination of linear phases, which is flicker-free by construction for a single trap per sub-plane of the SLM~\cite{kim2016}. For arbitrary geometries the holograms are computed by iterative algorithms of the Gerchberg--Saxton family. Kim, Kim, Lee and Ahn computed every frame of a rearrangement movie by a weighted Gerchberg--Saxton iteration that starts from the phase of the previous frame, which they call phase induction, and which keeps the phases of the traps continuous from frame to frame~\cite{kimoe}. Knottnerus and co-workers instead compute two endpoint holograms and interpolate the position, the amplitude and the optical phase of every trap linearly between them, which they call linear phase interpolation~\cite{lpi}.

The two published sequences leave a different part of the problem open. Linear phase interpolation never optimizes the intermediate frames, so the intensities of the traps in those frames are uneven, and the transient between two frames can dip. Per-frame Gerchberg--Saxton optimizes every frame, and the induction makes the phase of each trap start from its value in the previous frame, but nothing prescribes where the phase of a trap goes: the optimization is free to move it, and the only statement about its size is that the similarity of successive phases, a change small compared with one radian, is empirically guaranteed by the induction~\cite{kimoe}. For the number of frames, Kim et al.\ give a rule (enough frames that the step of a trap is at most the waist of the trap) and an empirical Gaussian fit of the transient minimum against the step; neither published method gives a closed form that turns a stated transient level into a frame count.

We propose the flicker-bounded phase-continuous sequence. Each frame is a weighted Gerchberg--Saxton step whose trap phases are locked to a smooth prescribed trajectory, started from the hologram and the weights of the previous frame, and the number of frames follows from a closed form for the intensity of a trap that is read half way between two frames. It would be wrong to claim that no published work optimizes the intermediate frames of a trap trajectory: Kim et al.~\cite{kimoe} compute every frame. The claim made here is narrower and is tested against that method: the phase trajectory of the traps is prescribed and locked, and the number of frames has a derived bound.

The contributions are the following.
\begin{enumerate}
\item The sequence itself: a hard lock of the trap phases to a minimal-jerk trajectory inside a warm-started weighted iteration, with the frame-count equation that bounds the pattern shift between frames (Section~\ref{sec:flicker-method}).
\item The theory that explains the bound: a closed form for the intensity of a shifted pattern between two frames, a bound on how far the lock lets the achieved trap phase move, and the conditions under which the closed form is only an upper estimate (Section~\ref{sec:theory}).
\item A controlled comparison with both published sequences, linear phase interpolation and per-frame Gerchberg--Saxton with phase induction at 5 and at 10 iterations per frame, run by the same code on the same trajectories, on the acceptance task, on long lattice moves and on the sort of a $6\times6$ array (Sections~\ref{sec:design} and~\ref{sec:results}).
\item An ablation that separates the lock, the per-frame optimization, the iteration count and the time profile (Section~\ref{sec:ablation}).
\item A statement of the trade that the measurements support, including the measures on which the published per-frame method is better (Section~\ref{sec:discussion}).
\end{enumerate}

The main result is a trade that depends on the time profile of the trap positions. With its default minimal-jerk profile and \parFrames{} frames the sequence keeps the lowest moving trap at \accBTransMean{} of its mean against \accLpiTransMean{} for linear phase interpolation, but the per-frame Gerchberg--Saxton sequence of Kim et al.\ keeps it at \accKimFiveTransMean{}. With the constant-velocity profile that the baselines use, in a comparison added after the results (Section~\ref{sec:design}), the same sequence reaches \ablTLinearMean{} and is read as better than all three baselines on the dip. In both profiles it has the smaller largest phase step, \stepRatio{} times smaller than that of the per-frame method at the default profile, and a per-frame intensity spread that is lower by a factor of \sigmaRatio{}. What it costs is a time per sequence of \secRatio{} times that of the per-frame method at its own iteration count, and, with the default profile on the longest moves, seeds on which the prescribed phase path cannot be followed and a sort on which the per-frame method is higher on every seed.

Section~\ref{sec:related} places the sequence among the published methods and Section~\ref{sec:problem} fixes the model, the blend of two frames and the metrics. Section~\ref{sec:flicker-method} defines the sequence with its algorithm, parameters and cost, and Section~\ref{sec:theory} gives the propositions. Sections~\ref{sec:design} to~\ref{sec:ablation} give the experiments, Section~\ref{sec:discussion} reads them against the problem, and Section~\ref{sec:limits} states what the measurements do not show.
